% Generado por scripts/70_tablas_figuras.R a partir de resultados-v1. No editar a mano.
\begin{table}[htbp]
\centering
\small
\caption{Tabla 6.8—Efecto de los proyectos sobre la conversión de bosque en sus polígonos}
\label{tab:tabla_6_8}
\begin{threeparttable}
\begin{tabular}{@{}lccccc@{}}
\toprule
 & (1) & (2) & (3) & (4) & (5) \\
 & ATT & Error estándar & IC 95 \% & Relativo & Media previa \\
\midrule
\multicolumn{6}{@{}l}{\textit{Panel A. Efecto principal}} \\
Pérdida de bosque natural (ha) & $-$0,098 & (0,351) & [$-$0,79; 0,59] & $-$19 \% & 0,518 \\
Tasa de pérdida (\% del bosque) & $-$0,024 & (0,124) & [$-$0,27; 0,22] & $-$14 \% & 0,163 \\
\addlinespace
\multicolumn{6}{@{}l}{\textit{Panel B. Por tipo de bosque de origen}} \\
Bosque de tierra firme & $-$0,075 & (0,350) & [$-$0,76; 0,61] & $-$17 \% & 0,434 \\
Bosque inundable & $-$0,022 & (0,046) & [$-$0,11; 0,07] & $-$26 \% & 0,084 \\
\addlinespace
\multicolumn{6}{@{}l}{\textit{Panel C. Por uso de destino}} \\
Agropecuario & $-$0,081 & (0,340) & [$-$0,75; 0,58] & $-$16 \% & 0,503 \\
Vegetación secundaria (rastrojo) & 0,013 & (0,258) & [$-$0,49; 0,52] & 11 \% & 0,116 \\
Minería & $-$0,008 & (0,015) & [$-$0,04; 0,02] & $-$168 \% & 0,005 \\
Otros usos antrópicos & $-$0,010 & (0,049) & [$-$0,11; 0,09] & $-$101 \% & 0,010 \\
\addlinespace
\multicolumn{6}{@{}l}{\textit{Panel D. Desplazamiento}} \\
Celdas de la zona de amortiguación & $-$0,194 & (0,163) & [$-$0,51; 0,13] & $-$21 \% & 0,927 \\
\addlinespace
\multicolumn{6}{@{}l}{\textit{Panel E. Sensibilidad al estimador}} \\
Doblemente robusto, con covariables & $-$0,160 & (2,856) & [$-$5,76; 5,44] & $-$31 \% & 0,518 \\
Doblemente robusto, sin covariables & $-$0,066 & (0,171) & [$-$0,40; 0,27] & $-$13 \% & 0,518 \\
\bottomrule
\end{tabular}
\begin{tablenotes}[flushleft]
\footnotesize
\item \textit{Notas:} Unidad: celda de 2 × 2 km (400 ha); ATT y media previa en hectáreas por celda y año, salvo la tasa (puntos porcentuales). Celdas tratadas: las que caen dentro de los polígonos de 25 proyectos de deforestación evitada (Karnik et al., 2025); controles emparejados por bosque en 2000, área inundable, elevación, pendiente, deforestación previa y presión de deforestación en el entorno. Estimador de Callaway y Sant’Anna por regresión sobre la muestra emparejada (17.064 celdas tratadas y 31.674 de control), con errores agrupados por proyecto. Panel D: celdas de la zona de amortiguación frente a controles (14.908 y 29.405 celdas). El relativo divide el ATT por la media previa; en minería y otros usos la media es casi cero y el porcentaje no es informativo. Panel E: el estimador doblemente robusto con covariables es inestable por la separación casi completa en el modelo de propensión.
\item \textit{Fuente:} MapBiomas Colombia, Colección 3; Karnik et al. (2025); cálculos propios.
\end{tablenotes}
\end{threeparttable}
\end{table}
